\begin{abstract}
\noindent
We report Stage~1 and Stage~2 of a controlled study of \optCoupled, a partner-weight-aware Newton--Schulz preconditioner whose stage-1 inner step orthogonalises the bilinear products $W_Q W_K^{\top}$, $W_V W_O$ and $W_{\text{up}} W_{\text{down}}$ rather than each factor in isolation. Stage~1 reproduces the LLaMA-60M baseline (rung A0; 75/75 cells, 5 LRs $\times$ 5 seeds per optimizer) and Stage~2 walks a six-rung dense architecture ladder B--H at the same 60M Common-Size shape (3$\times$5 LRs $\times$ 3 seeds; 270/270 cells fully synced, 88 diverged). All three optimizers (\optAdamW, \optMuon, \optCoupled) outperform a well-tuned \optAdamW{} baseline by $0.12$--$0.29\nats{}$ of final validation loss, a margin well outside seed noise. Within the \optMuon{} family, paired-bootstrap 95\% CIs at each arm's own best LR show \optCoupled{} beating \optMuon{} by a small but consistent margin on every rung that retains full RoPE (A0: $-0.024$\nats, B: $-0.008$\nats, G: $-0.011$\nats, H: $-0.014$\nats, all CIs excluding zero), and \emph{regressing} on the two learned-positional rungs C ($+0.066$\nats) and D ($+0.055$\nats), where the multi-head Q--K branch silently falls back to flat-2D coupling and \optCoupled{}'s best LR drops 10$\times$ below \optMuon{}'s. The Karpathy vanilla rung E shows a tie within CI. The result is a clean architecture$\times$optimizer interaction signal that ties the coupled gain to the RoPE-equipped attention block, consistent with the algorithm's bilinear-product motivation. Stage~3 (sparse MoE rungs I--N) has since landed on rigs A/B but its interpretation is out of scope here.
\end{abstract}
